% Part: many-valued-logic
% Chapter: three-valued-logics
% Section: multiple-designation

\documentclass[../../../include/open-logic-section]{subfiles}

\begin{document}

\olfileid{mvl}{thr}{mul}

\olsection{Milih nilai saliyane mung $\True$}

Nganti saiki, kabeh logika kang wis kita deleng nduweni himpunan
nilai kabeneran pinilih $V^+ = \{\True\}$: sawijining pratelan
dianggep bener yen lan mung yen nilai kabenerane~$\True$.
Nanging sawijining pratelan uga bisa dianggep bener anggere
nilainé dudu~$\False$. Kita bisa oleh logika mangkono kanthi
netepake ing matriks, umpamane, $V^+ = \{\True, \Undef\}$.

\begin{defn}
\emph{Logika paradoks}~$\LogLP$ ditetepake nganggo matriks iki:
\begin{enumerate}
  \item Perangan tanpa tetapan saka basa proposisional baku
  $\Lang L_0$ kang ngemot $\lnot$, $\land$, $\lor$, $\lif$.
  \item Himpunan nilai kabeneran $V = \{\True, \Undef, \False\}$.
  \item $\True$ lan $\Undef$ iku nilai pinilih, yaiku $V^+ = \{\True, \Undef\}$.
  \item Fungsi kabenerane padha karo logika Kleene kang kuwat.
\end{enumerate}
\end{defn}

\begin{defn}
Halld\'en ngenalake \emph{logika tanpa makna}~$\LogHal$ kanthi matriks iki:
\begin{enumerate}
  \item Perangan tanpa tetapan saka basa proposisional baku
  $\Lang L_0$ kang ngemot $\lnot$, $\land$, $\lor$, $\lif$
  lan konektif $1$-argumen~$+$.
  \item Himpunan nilai kabeneran $V = \{\True, \Undef, \False\}$.
  \item $\True$ lan $\Undef$ iku nilai pinilih, yaiku $V^+ = \{\True, \Undef\}$.
  \item Fungsi kabenerane padha karo logika Kleene kang ringkih,
  ditambahi operator ``tanpa makna'' iki:
  \begin{center}
    \begin{tabular}{c|c} 
    $\tf{+}$ & \\ 
    \hline  
    $\True$ & $\False$ \\ 
    $\Undef$ & $\True$ \\
    $\False$ & $\False$ 
    \end{tabular}
  \end{center}
\end{enumerate}
\end{defn}

Beda karo logika-logika Kleene kang nduweni tabel kabeneran
padha, logika loro iki \emph{nduweni} tautologi.

\begin{prop}\ollabel{prop:LP-taut-CL} Tautologi ing $\LogLP$
  padha karo tautologi ing logika proposisional klasik.
\end{prop}

\begin{proof}
  Miturut \olref[syn][sub]{prop:mvl-cl}, yen $\Entails[\LogLP] !A$
  mula $\Entails[\LogCL] !A$. Kanggo mbuktekake arah kosok baline,
  kita tuduhake: yen ana !!a{valuation}~$\pAssign v\colon \PVar \to
  \{\False, \True, \Undef\}$ kang nyukupi $\pValue v(!A)[\LogKs]
  = \False$, mula ana !!a{valuation}~$\pAssign {v'}\colon \PVar \to
  \{\False, \True\}$ kang nyukupi $\pValue {v'}(!A)[\LogCL] =
  \False$. Iki mbuktekake asil kanggo $\LogLP$, amarga $\LogKs$
  lan $\LogLP$ nduweni fungsi kabeneran kang padha, lan mung
  $\False$ kang ora dadi nilai pinilih ing $\LogLP$ (iki mung
  bedane $\LogLP$ karo~$\LogKs$). Dadi, yen $\Entails/[\LogLP]
  !A$, ana !!{valuation}~$\pAssign v$ kang nyukupi $\pValue
  v(!A)[\LogLP] = \pValue v(!A)[\LogKs] = \False$. Miturut klaim
  kang bakal dibuktekake, $\pValue{v'}(!A)[\LogCL] = \False$,
  yaiku $\Entails/[\LogCL] !A$.

Kanggo mbuktekake klaim kasebut, luwih dhisik kita tetepake
  $\pAssign {v'}$ mangkene:
  \[
    \pAssign {v'}(p) =
    \begin{cases}
    \True & \text{yen } \pAssign {v}(p) \in \{\True, \Undef\}\\
    \False & \text{yen ora}
  \end{cases}
  \]
  Saiki kita buktèkake kanthi induksi tumrap $!A$ yen (a)~yen
  $\pValue v(!A)[\LogKs] = \False$, mula $\pValue {v'}(!A)[\LogCL]
  = \False$; lan (b)~yen $\pValue v(!A)[\LogKs] = \True$, mula
  $\pValue {v'}(!A)[\LogCL] = \True$.
  \begin{enumerate}
    \item Dhasar induksi: $!A \ident p$. Yen nilai variabel p salah utawa
    bener, miturut \olref[syn][val]{defn:pValue},
    $\pValue v(!A)[\LogKs] = \pAssign v(p) = \pValue {v'}(!A)[\LogCL]$;
    iki negesake (a) lan~(b). Yen nilai variabel p ora katetepake,
    prasyarat (a) lan (b) loro-lorone ora kawujud.

Kanggo langkah induksi, gatekna kasus-kasus iki:
    \item $!A \ident \lnot !B$.
    \begin{enumerate}
      \item Upamakna $\pValue v(\lnot!B)[\LogKs] = \False$.
      Miturut definisi $\tf{\lnot}[\LogKs]$, kita oleh
      $\pValue v(!B)[\LogKs] = \True$. Miturut hipotesis induksi,
      kasus~(b), $\pValue {v'}(!B)[\LogCL] = \True$; mula
      $\pValue {v'}(\lnot !B)[\LogCL] = \False$.
      \item Upamakna $\pValue v(\lnot!B)[\LogKs] = \True$.
      Miturut definisi $\tf{\lnot}[\LogKs]$, kita oleh
      $\pValue v(!B)[\LogKs] = \False$. Miturut hipotesis induksi,
      kasus~(a), $\pValue {v'}(!B)[\LogCL] = \False$; mula
      $\pValue {v'}(\lnot !B)[\LogCL] = \True$.
    \end{enumerate}

\item $!A \ident (!B \land !C)$.
    \begin{enumerate}
      \item Upamakna $\pValue v(!B \land !C)[\LogKs] = \False$.
      Miturut definisi $\tf{\land}[\LogKs]$, salah siji saka
      $\pValue v(!B)[\LogKs] = \False$ utawa
      $\pValue v(!C)[\LogKs] = \False$ kudu lumaku. Miturut
      hipotesis induksi, kasus~(a), salah siji saka
      $\pValue {v'}(!B)[\LogCL] = \False$ utawa
      $\pValue {v'}(!C)[\LogCL] = \False$ lumaku; mula
      $\pValue {v'}(!B \land !C)[\LogCL] = \False$.
      \item Upamakna $\pValue v(!B \land !C)[\LogKs] = \True$.
      Miturut definisi $\tf{\land}[\LogKs]$,
      $\pValue v(!B)[\LogKs] = \True$ lan
      $\pValue v(!C)[\LogKs] = \True$. Miturut hipotesis induksi,
      kasus~(b), $\pValue {v'}(!B)[\LogCL] = \True$ lan
      $\pValue {v'}(!C)[\LogCL] = \True$; mula
      $\pValue {v'}(!B \land !C)[\LogCL] = \True$.
    \end{enumerate}
  \end{enumerate}

Rong kasus liyane padha carane lan dadi latihan. Utawa,
  bukti ing ndhuwur wis mbuktekake asil kanggo kabeh !!{formula}s
  kang mung ngemot $\lnot$ lan~$\land$. Sabanjure kita bisa
  nggunakake kasunyatan yen ing $\LogKs$ lan $\LogCL$, kanggo
  saben $\pAssign v$, $\pValue v(!B \lor !C) = \pValue
  v(\lnot(\lnot!B \land \lnot !C))$ lan $\pValue v(!B \lif !C)
  = \pValue v(\lnot(!B \land \lnot !C))$.
\end{proof}

\begin{prob}
  Rampungna bukti \olref[mvl][thr][mul]{prop:LP-taut-CL}, yaiku
  buktèkna (a) lan~(b) kanggo kasus $!A \ident (!B \lor !C)$
  lan $!A \ident (!B \lif !C)$.
\end{prob}

\begin{prob}
Buktèkna yen saben tautologi klasik uga tautologi ing~$\LogHal$.
\end{prob}

Sanajan tautologie padha karo logika klasik ing perangan basa
tanpa operator tambahan, sesambungan konsekuensine beda. Umpamane,
$\LogLP$ iku \emph{parakonsisten}, amarga $\lnot p, p \Entails/ q$;
mula prinsip eksplosi $\lnot !A, !A \Entails !B$ ora lumaku kanthi
umum. (Prinsip iki lumaku kanggo sawetara $!A$ lan $!B$, umpamane
yen $!B$~iku tautologi.)

\begin{prob}
  Sesambungan endi ing ngisor iki kang lumaku ing (a)~$\LogLP$
  lan (b)~$\LogHal$? Wenehana tabel kabeneran kanggo saben
  sesambungan.
  \begin{enumerate}
    \item $p, p \lif q \Entails q$
    \item $\lnot q, p \lif q \Entails \lnot p$
    \item $p \lor q, \lnot p \Entails q$
    \item $\lnot p, p \Entails q$
    \item $p \Entails p \lor q$
    \item $p \lif q, q\lif r \Entails p \lif r$
  \end{enumerate}
\end{prob}

Kepriye yen $\Undef$ uga dadi nilai pinilih ing $\LogLuk[3]$?

\begin{defn}
  \emph{R-Mingle telung nilai}~$\LogRM[3]$ ditetepake nganggo matriks iki:
  \begin{enumerate}
    \item Basa proposisional baku $\Lang L_0$ kang ngemot
    $\lfalse$, $\lnot$, $\land$, $\lor$, $\lif$.
    \item Himpunan nilai kabeneran $V = \{\True, \Undef, \False\}$.
    \item $\True$ lan $\Undef$ iku nilai pinilih, yaiku $V^+ = \{\True, \Undef\}$.
    \item Fungsi kabenerane padha karo logika \L ukasiewicz~$\LogLuk[3]$.
  \end{enumerate}
\end{defn}

\begin{prob}
  Sesambungan endi ing ngisor iki kang lumaku ing $\LogRM[3]$?
  \begin{enumerate}
    \item $p, p \lif q \Entails q$
    \item $p \lor q, \lnot p \Entails q$
    \item $\lnot p, p \Entails q$
    \item $p \Entails p \lor q$
  \end{enumerate}
\end{prob}

Tabel kabeneran kang beda kadhangkala bisa ngasilake logika
(sesambungan entailment) kang padha mung kanthi ngganti nilai
pinilih. Umpamane, iki kedadeyan yen ing logika G\"odel kita
milih $V^+ = \{\True, \Undef\}$ tinimbang~$\{\True\}$.

\begin{prop}\ollabel{prop:gl-udes}
  Matriks kanthi $V = \{\False, \Undef,\True\}$, $V^+=\{\True,
  \Undef\}$, lan fungsi kabeneran logika G\"odel $3$ nilai
  netepake logika klasik.
\end{prop}

\begin{proof}
  Latihan.
\end{proof}

\begin{prob}
  Buktèkna \olref[mvl][thr][mul]{prop:gl-udes} kanthi nuduhake yen
  kanggo logika~$\Log L$ kang ditetepake kaya logika G\"odel nanging
  nganggo $V^+=\{\True,\Undef\}$, yen $\Gamma \Entails/[\Log L]
  !B$ mula $\Gamma \Entails/[\LogCL] !B$. Gunakna gagasan saka
  \olref[mvl][thr][mul]{prop:LP-taut-CL}, nanging tinimbang mbuktekake
  sipat (a) lan~(b), tuduhna yen $\pValue v(!A)[\LogGod] = \False$
  yen lan mung yen $\pValue {v'}(!A)[\LogCL] = \False$ (lan mula
  $\pValue v(!A)[\LogGod] \in \{\True,\Undef\}$ yen lan mung yen
  $\pValue {v'}(!A)[\LogCL] = \True$). Terangna ngapa iki
  mbuktekake proposisi kasebut.
\end{prob}

\end{document}
